% Einbindung nach 02-erfuellungsrelation.tex: B48FormulaInTree steht bereit.
% Die Positionsnotation ist in Band 27 lokal; providecommand ermoeglicht
% auch die eigenstaendige Uebersetzung von Band 48.
\providecommand{\TreePositions}[2]{\operatorname{Pos}_{#1}(#2)}
\providecommand{\AddressEdges}[1]{E_{#1}^{\mathrm{adr}}}
\providecommand{\AddressLeft}[1]{\lambda_{#1}^{\mathrm{adr}}}
\providecommand{\AddressRight}[1]{\rho_{#1}^{\mathrm{adr}}}

\subsection{Die endliche Baumstruktur der Formelcodes}

Die Formelmenge ist eine durch ihre Konstruktoren abgeschlossene
Teilmenge des Baumtraegers aus Band~27. Ihr Induktionsprinzip
\FormulaRefAuto{B48FormulaInduction}[thm] beruht auf dieser kleinsten
abgeschlossenen Menge. Fuer jeden einzelnen Code steht zusaetzlich
die endliche Graphstruktur seiner Vorkommen zur Verfuegung.

\FormulaDefDeltaK[Positionsdaten eines Formelcodes]
{\begin{aligned}
 P_p&\coloneq\TreePositions{\Sigma}{p},&
 E_p&\coloneq\AddressEdges{P_p},\\
 \lambda_p&\coloneq\AddressLeft{P_p},&
 \rho_p&\coloneq\AddressRight{P_p},\\
 \varepsilon_2&\coloneq\EmptyWord{\{0,1\}}.
\end{aligned}}
{B48FormulaPositionData}{\DeltaRow{Baumcode}{p\in\mathcal T}}

Die rechte Seite verwendet die bereits definierten Positionsmengen
und Adressgraphen aus \FormulaRefAuto{TreePositionSetDef}[def] und
\FormulaRefAuto{BinaryAddressGraphDef}[def]. Die neuen Kurzzeichen
erleichtern nur den Zugriff auf diese Daten.

\FormulaThmDeltaK[Formelcodes tragen endliche geordnete volle Binaerbaeume]
{\begin{aligned}
 p\in\mathcal F\vdash{}&\Finite{P_p}\\[-2pt]
 &\land\OrderedFullBinaryTreeStruct{
     P_p,E_p,\varepsilon_2,\lambda_p,\rho_p}
\end{aligned}}
{B48FormulaPositionTree}{\DeltaRow{Formelcode}{p}}
\begin{tabproofwide}
 \proofstepwidestar[1]{p\in\mathcal F}{\rA}
 \proofstepwidestar[1]{p\in\mathcal T}{%
  \FormulaRefAuto{B48FormulaInTree}[thm]{1}}
 \proofstepwidestar[1]{p\in\BracketTreeSet{\Sigma}}{%
  \FormulaRefAuto{B48CodeCarrier}[def]{2}}
 \proofstepwidestar[1]{\Finite{\TreePositions{\Sigma}{p}}}{%
  \FormulaRefAuto{BracketCodePositionSetFinite}[thm]{3}}
 \proofstepwidestar[1]{\begin{aligned}[t]
  &\mathsf{GVollBinBaum}\!\bigl(\TreePositions{\Sigma}{p},\\[-2pt]
  &\quad\AddressEdges{\TreePositions{\Sigma}{p}},
    \EmptyWord{\{0,1\}},\\[-2pt]
  &\quad\AddressLeft{\TreePositions{\Sigma}{p}},
    \AddressRight{\TreePositions{\Sigma}{p}}\bigr)
 \end{aligned}}{%
  \FormulaRefAuto{BracketCodePositionGraphOrderedFullBinaryTree}[thm]{3}}
 \proofstepwidestar[1]{\begin{aligned}[t]
   &\Finite{P_p}\\[-2pt]
   &\land\OrderedFullBinaryTreeStruct{
      P_p,E_p,\varepsilon_2,\lambda_p,\rho_p}
 \end{aligned}}{%
  \FormulaRefAuto{B48FormulaPositionData}[def]{\rAI{4,5}}}
\end{tabproofwide}

Die Knoten sind Vorkommen im Code. Derselbe Teilcode kann daher an
mehreren verschiedenen Positionen erscheinen. Der volle Binaerbaum
enthaelt auch die bereits eingefuehrten Hilfsmarken; ein einzelner
logischer Junktor muss deshalb nicht einem einzelnen Binaerknoten
entsprechen. Die Aussagen ueber Eltern, Kinder und Wurzelwege aus
Band~26 gelten fuer diesen Codegraphen. Auf seinen Teilmengen sind
die Saetze ueber endliche Mengen aus Band~20 anwendbar. Damit stehen
sowohl die Formelinduktion ueber Konstruktoren als auch die endliche
Menge ihrer jeweiligen Vorkommen auf einer bereits bewiesenen Grundlage.
